\section{zLLM Post-Quantum Commitments}
\label{sec:zllm-pq}

zLLM is \Zoo{}'s training-free experience-library adaptation system:
a base model is frozen, and an experience library curated by the \DAO{}
is injected at inference. The system needs commitments at three
points: the experience-library state, the governance vote that
admitted each experience, and the inference receipts that record
which experiences were used.

In 3.0 each of these three commitments uses post-quantum primitives.

\subsection{Experience library commitment}

The experience library is a Merkle accumulator over individual
experience entries (a tuple of prompt context, expected behaviour,
provenance, and metadata). Pre-3.0 the accumulator used SHA-256 over
classical Merkle paths and was signed by the curating multisig with
ECDSA. In 3.0:

\begin{itemize}[leftmargin=1.2em]
  \item Hash function: SHA3-256 (Keccak-f) for collision resistance under
        Grover-bounded adversaries; the $128$-bit post-quantum security
        target requires the underlying hash to provide $256$-bit
        classical resistance.
  \item Accumulator: a sparse Merkle tree (still classical structure;
        PQ resistance comes from the hash and the signing key).
  \item Signing key on the accumulator root: threshold
        $\MLDSA{}\text{-}65$ over the curators set, $t$-of-$n$ where
        $t \geq \lceil 2n/3 \rceil$.
\end{itemize}

The experience library commit transaction posts the new root and the
threshold $\MLDSA{}$ signature. Verifiers re-derive the threshold
public key from the on-chain curators set and check the signature.

\subsection{Governance vote commitments}

zLLM experience admission is a \DAO{} vote. Each ballot is a
$\MLDSA{}\text{-}65$ signature over the experience hash, the voter's
weight, and the vote direction. Tally is a deterministic Merkle root
over the ballots, signed by the tally proposer with $\MLDSA{}$.
Optional \FHE{} mode (per \cite{zoo-2025-securities-and-dao}) keeps
ballot direction private until the tally is opened, but the public
result is still authenticated by $\MLDSA{}$.

A vote receipt for an admitted experience contains:

\begin{enumerate}[leftmargin=1.4em]
  \item The experience hash.
  \item The governance proposal hash.
  \item The tally Merkle root.
  \item The \MLDSA{} signature on the tally root.
  \item The block height of the tally on the L2.
  \item The reference \Quasar{}Cert from \Lux{} L1.
\end{enumerate}

A verifier can reconstruct the entire admission chain from these
six fields.

\subsection{Inference receipts}

When a zLLM inference uses the experience library, the inference
receipt records which experiences were retrieved. The receipt schema
in 3.0:

\begin{lstlisting}[language=C++,caption={3.0 inference receipt schema}]
struct InferenceReceipt3 {
    Hash      base_model_id;          // pinned weights digest
    Hash      experience_root;        // accumulator root at retrieval
    Hash[]    experience_paths;       // Merkle paths for retrieved entries
    Hash      prompt_digest;          // hash of input (or FHE handle)
    Hash      output_digest;          // hash of output (or FHE handle)
    TeeQuote  tee_quote;              // attestation of execution env
    MlDsa::Sig signing_proof;          // worker's ML-DSA-65 signature
    Hash      a_chain_anchor;         // commitment to A-Chain (LP-134)
};
\end{lstlisting}

The receipt establishes:

\begin{itemize}[leftmargin=1.2em]
  \item which model produced the output (\texttt{base\_model\_id}),
  \item which experiences shaped it (\texttt{experience\_paths}),
  \item that the execution happened inside a TEE-attested container
        (\texttt{tee\_quote}),
  \item that the worker authenticated under a post-quantum signing
        key (\texttt{signing\_proof}),
  \item that the receipt is anchored to A-Chain (LP-134) for
        cross-chain verifiability.
\end{itemize}

\subsection{Attribution graph}

zLLM rewards flow back to upstream contributors: experience curators,
training-data donors, evaluators, and reviewers. The attribution
graph is itself a series of commitments. In 3.0 every node in the
attribution graph is signed with the contributor's $\MLDSA{}$ key,
the edge between contributor and experience is a Merkle inclusion
proof under the experience library root, and the disbursement
transaction is a triple-cert-anchored \Lux{} L1 transaction
\cite{lp-020}.

This means a long-tail contributor whose data shaped a Zen model in
2024 still receives correctly-attributed disbursements years later
under post-quantum verification --- the historical record is not
invalidated by Q-day.

\subsection{Migration of pre-3.0 commitments}

Experience-library entries created before 3.0 carry classical
commitments. The migration plan:

\begin{enumerate}[leftmargin=1.4em]
  \item \textbf{Re-sign campaign.} Curators re-sign the legacy
        accumulator root with the new threshold $\MLDSA{}\text{-}65$
        key. The on-chain commitment now carries both the legacy
        ECDSA and the new $\MLDSA{}$ signatures.
  \item \textbf{Sunset window.} For 90 days post-3.0 spec freeze,
        verifiers accept either signature.
  \item \textbf{PQ-only enforcement.} After the sunset window,
        only $\MLDSA{}$ signatures are accepted. Legacy ECDSA-only
        commitments must be re-signed before that date or be
        considered un-attestable.
\end{enumerate}

This is the same pattern adopted by the \DAO{} key migration
(\S\ref{sec:dao-migration}).
